% CODEX: wstaw bezpośrednio przed \section{Przestrzeń euklidesowa i afiniczna}
% w 29a-geometria-kleina.tex. Nie powtarzaj istniejącej teorii działań,
% zwartych stabilizatorów ani ilorazów przez grupy dyskretne.
\section{Struktury lokalnie modelowane na przestrzeni jednorodnej}
\label{sec:thur-dop-struktury-gx}
Geometria Kleina opisuje pojedynczą przestrzeń $X$ z grupą
symetrii $G$. Natomiast rozmaitość może tylko \emph{lokalnie}
wyglądać jak $X$: jej topologia nie musi dopuszczać
jednej globalnej mapy do tego modelu. Formalizujemy to
analogicznie do atlasu rozmaitości gładkiej, lecz ograniczamy
przekształcenia przejścia.

\begin{definicja}[Atlas struktury $(G,X)$]
\label{def:thur-dop-gx-atlas}
Niech $X$ będzie spójną gładką rozmaitością, na której grupa Liego
$G$ działa skutecznie przez dyfeomorfizmy. \emph{Strukturą $(G,X)$}
na rozmaitości $M$ nazywamy atlas odwzorowań
$\varphi_i:U_i\to V_i\subset X$, będących dyfeomorfizmami
na otwarte podzbiory modelu, taki że na każdej spójnej składowej
przecięcia $U_i\cap U_j$ odwzorowanie przejścia
$\varphi_j\circ\varphi_i^{-1}$ jest ograniczeniem
pewnego elementu $g_{ji}\in G$. Zwykle utożsamiamy atlasy
zgodne w tym sensie.
\end{definicja}

Właściwą klasą lokalnych przekształceń jest pseudogrupa
utworzona z ograniczeń elementów $G$ i sklejania takich
ograniczeń. Zwykły atlas gładki odpowiada znacznie większej
pseudogrupie wszystkich lokalnych dyfeomorfizmów;
w atlasie geometrycznym przejścia muszą zachowywać
wybraną strukturę modelu.

\begin{twierdzenie}[Rozwinięcie i reprezentacja holonomii]
\label{thm:thur-dop-dev-hol}
Załóżmy, że $X$ jest spójna i jednospójna, a $G$ działa
skutecznie przez izometrie pewnej metryki riemannowskiej
na $X$. Dla spójnej rozmaitości $M$ ze strukturą $(G,X)$
istnieją lokalny dyfeomorfizm
$\mathrm{dev}:\widetilde M\to X$ i homomorfizm
$\rho:\pi_1(M)\to G$, spełniające
\begin{equation}\label{eq:thur-dop-dev-equiv}
 \mathrm{dev}(\gamma\cdot\widetilde p)
 =\rho(\gamma)\cdot\mathrm{dev}(\widetilde p).
\end{equation}
Po wyborze punktów bazowych i identyfikacji przekształceń
nakrycia para $(\mathrm{dev},\rho)$ jest wyznaczona z dokładnością
do jednoczesnego złożenia $\mathrm{dev}$ z elementem $G$
i sprzężenia $\rho$ tym elementem.
\end{twierdzenie}
\begin{proof}
Wybierzmy punkt bazowy i jedną mapę atlasu w jego otoczeniu.
Przedłużamy tę mapę wzdłuż dowolnej drogi, przechodząc kolejno
przez mapy atlasu i składając je z odpowiednimi elementami $G$.
Na skończonym podziale parametru drogi wystarczy skończenie
wiele map. Przekształcenia przejścia po dwóch równoważnych
lokalnych przedłużeniach zgadzają się na niepustym zbiorze
otwartym. Ponieważ izometria spójnej rozmaitości jest
wyznaczona przez swoje ograniczenie do takiego zbioru,
przedłużenie jest jednoznaczne. Homotopia dróg z ustalonymi
końcami, podzielona na dostatecznie małe prostokąty,
nie zmienia wyniku. Na nakryciu uniwersalnym otrzymujemy
więc jednoznacznie odwzorowanie $\mathrm{dev}$, lokalnie
równe jednej z map atlasu złożonej z elementem $G$.

Zastosowanie tego przepisu do pętli opartej w punkcie
bazowym zmienia początkową mapę o pewien element
$\rho(\gamma)\in G$. Składanie pętli odpowiada składaniu
tych elementów, co daje homomorfizm; różnica wyników
przy zmianie punktu podniesionego to dokładnie
równość~\eqref{eq:thur-dop-dev-equiv}. Inny wybór
początkowej mapy zmienia całą konstrukcję przez jedną
symetrię modelu.
\end{proof}

\begin{uwaga}[Dwie różne holonomie]
\label{uw:thur-dop-dwie-holonomie}
Reprezentacja $\rho:\pi_1(M)\to G$ to
\emph{holonomia struktury $(G,X)$}. Nie należy jej
utożsamiać z grupą holonomii koneksji Leviego-Civity
z Definicji~\ref{def:grupa-holonomii}. Na płaskim torusie
$T^n=\R^n/\Z^n$ pierwsza holonomia wysyła element
$m\in\Z^n$ na przesunięcie $x\mapsto x+m$, więc jest
nietrywialna; druga jest trywialna, bo pola współrzędne
schodzą do globalnej ramy równoległej.
\end{uwaga}

\begin{stwierdzenie}[Zupełność a globalny iloraz]
\label{prop:thur-dop-dev-kompletna}
Jeśli $X$ jest jednospójna i zupełna jako rozmaitość
riemannowska, a atlas $(G,X)$ na spójnej $M$ wyznacza
zupełną metrykę lokalnie izometryczną z $X$, to
$\mathrm{dev}:\widetilde M\to X$ jest izometrią.
Wówczas $M$ jest ilorazem $X$ przez grupę
$\rho(\pi_1(M))$. W szczególności dotyczy to
zwartych rozmaitości z taką strukturą.
\end{stwierdzenie}
\begin{proof}
Metryka z $X$ skleja się na $M$, ponieważ zmiany
współrzędnych są izometriami. Odwzorowanie rozwijające
jest lokalną izometrią, a nakrycie uniwersalne z podniesioną
metryką jest zupełne. Zupełna lokalna izometria
na spójną rozmaitość riemannowską jest nakryciem:
podnosząc geodezyjne wychodzące z punktu obrazu,
przedłużamy je dzięki zupełności i uzyskujemy lokalne
odwrotności na kulach normalnych, a potem pokrywamy nimi
przestrzeń docelową. Nakrycie jednospójnej $X$ jest
trywialne, zatem $\mathrm{dev}$ jest globalną izometrią.
Równość~\eqref{eq:thur-dop-dev-equiv} identyfikuje
działanie grupy nakrycia z działaniem $\rho(\pi_1(M))$.
Jeśli $M$ jest zwarta, jej każda metryka riemannowska
jest zupełna dzięki twierdzeniu Hopfa--Rinowa.
\end{proof}

\begin{przyklad}[Struktura niezupełna]
\label{prz:thur-dop-struktura-niezupelna}
Na otwartej kuli euklidesowej $B^n\subset\R^n$ wybierzmy
atlas otrzymany przez ograniczenie identyczności
z modelu $X=\R^n$ i grupy jego izometrii. Nakrycie
uniwersalne jest trywialne, $\mathrm{dev}:B^n\hookrightarrow\R^n$
jest włożeniem, a holonomia jest trywialna.
Metryka nie jest zupełna: odcinek zbliżający się do brzegu
ma skończoną długość. Dlatego nie dostajemy
ilorazu całego modelu $\R^n$.
\end{przyklad}
